\section*{Chapter \ref{ch:m2:credit-indices-and-tranches} --- Credit Indices and Tranches}

\begin{solution}{exo:m2:credit-indices-and-tranches:1}
The name's share is $1\text{ billion}/125 = \text{USD}~8$ million, on which you pay
$(100 - 40)\% $: USD~4.8 million. The index continues on 124 names with a notional of
USD~992 million, and the premium is paid on that.
\end{solution}

\begin{solution}{exo:m2:credit-indices-and-tranches:2}
0\% (the losses stop below 3\%), 50\% ($(5-3)/4$) and 100\% (above 7\%).
\end{solution}

\begin{solution}{exo:m2:credit-indices-and-tranches:3}
Its list reflects today's liquid investment-grade names, and everyone quotes and hedges
in it, so its bid--offer is tightest. The old series keeps its list, including names
that have been downgraded or become illiquid, and trades less and less, off the run.
\end{solution}

\begin{solution}{exo:m2:credit-indices-and-tranches:4}
No: 146.0 basis points. The wider name has the smaller annuity, 4.085 against 4.420,
so it counts for less in the average upfront.
\end{solution}

\begin{solution}{exo:m2:credit-indices-and-tranches:5}
Higher correlation makes defaults cluster: most of the time few or none happen, which
spares the equity, and occasionally many, which the equity loses entirely anyway. Its
expected loss falls as correlation rises, so the \omterm{def:m2:credit-default-swaps:cds}{protection seller} gains.
\end{solution}

\begin{solution}{exo:m2:credit-indices-and-tranches:6}
It exercises at expiry if the index spread is then above 80 basis points: it buys
protection on the index at the strike's terms, which is worth about (spread minus 80
basis points) times the index annuity, and it can keep the protection or sell it at
the market.
\end{solution}

\begin{solution}{exo:m2:credit-indices-and-tranches:7}
At $\rho = 0.3$: 59.57\%, 24.09\%, 7.828\% and 0.224\% of the \omterm{def:m2:credit-indices-and-tranches:tranche}{tranches}; times their
widths 3, 4, 8 and 85\%, they add up to 3.568\%, the portfolio's $5.95\% \times 60\%$.
The mezzanine's expected loss is 24.3\% at $\rho = 0.10$, and again at $\rho = 0.28$.
\end{solution}

\begin{solution}{exo:m2:credit-indices-and-tranches:8}
The index is compared with its \omterm{def:m2:credit-indices-and-tranches:intrinsic}{intrinsic spread}, not the simple average, which is
higher (here by 1.7 basis points). The trade costs half a bid--offer on 125 names and
the index, in and out; it is not riskless, because the skew can widen before it
closes, marking the position down and calling for margin, and the fund may not be able
to hold it until it closes. The default risk is hedged; the skew is not.
\end{solution}

\begin{solution}{pb:m2:credit-indices-and-tranches:1}
\textbf{1.} 75.3 and 55 basis points.
\textbf{2.} $-1.157\%$, an \omterm{def:m2:credit-indices-and-tranches:intrinsic}{intrinsic spread} of 73.6 basis points.
\textbf{3.} Upfronts, not spreads, are averaged, and names with wide spreads have small
annuities, so they weigh less.
\textbf{4.} Also 73.6: the first-order formula is very close.
\textbf{5.} $65.6 - 73.6 = -8$ basis points.
\textbf{6.} It receives 1.513\% on the index, USD~15.13 million, and pays 1.157\% on the
constituents, USD~11.57 million.
\textbf{7.} USD~3.55 million received: the value of the skew at the entry prices,
which the fund gives back if it closes the trade at the same prices.
\textbf{8.} It receives USD~4.8 million on the index and pays USD~4.8 million on its
single-name contract: nothing net.
\textbf{9.} USD~3.55 million.
\textbf{10.} A loss of USD~5.37 million.
\textbf{11.} USD~0.657 million on the constituents and 0.109 on the index, USD~0.77
million to enter; USD~1.53 million to enter and leave.
\textbf{12.} USD~2.02 million. The P\&L is about the change in skew times the annuity,
so the skew must move by more than $2 \times (1.5 + 0.25) = 3.5$ basis points.
\textbf{13.} Because a large participant keeps selling index protection, single names
are costly and slow to trade, and arbitrage capital is limited and must bear the
mark-to-market.
\textbf{14.} Variation margin on the index position, the use of capital and funding
limits, and possibly a forced exit at the worst moment.
\textbf{15.} Its sales move the index away from its constituents, and the skew traders
who take the other side profit when it has to stop selling or buy back; the larger the
position relative to the market's volume, the more its exit costs.
\textbf{16.} Unwinding means buying back protection equal to 10 to 15 days of the whole
market's volume, which would move the index sharply against it; everyone watching
the skew would anticipate it.
\textbf{17.} By the loss it can bear if the skew doubles before closing, by the
market's daily volume in the index and the names, and by the margin it can fund.
\textbf{18.} The \omterm{def:m2:credit-indices-and-tranches:index}{index series} goes off the run: liquidity moves to the new series, and
the old one's skew may close more slowly or be harder to trade.
\textbf{19.} Named result: \emph{the skew trade} makes USD~3.55 million gross and
USD~2.02 million net of costs when the skew closes from $-8$ to 0, and loses USD~5.37
million on paper if it first widens to $-20$.
\textbf{20.} Because closing it depends on the other side stopping and on the fund
surviving the widening until it does.
\end{solution}

\begin{solution}{iq:m2:credit-indices-and-tranches:1}
A \omterm{def:m2:credit-default-swaps:cds}{credit default swap} on 125 equally weighted, liquid North American investment-grade
names with a 100 basis point coupon, quoted as a spread. Every six months, on 20 March
and 20 September, a new series with an updated list is published and becomes the
on-the-run contract; the old series continues with its list.
\iqlookfor{constituents, coupon, the roll.}
\end{solution}

\begin{solution}{iq:m2:credit-indices-and-tranches:2}
The spread implied by the constituents' single-name prices: average their upfronts at
the index coupon and convert back. The index trades on its own supply and demand, is far
more liquid than the names, and arbitrage between them is costly and risky, so the two
can diverge, sometimes widely.
\iqlookfor{upfront averaging, and the limits to arbitrage.}
\end{solution}

\begin{solution}{iq:m2:credit-indices-and-tranches:3}
Name $i$ defaults if $\sqrt\rho Z + \sqrt{1-\rho}\varepsilon_i < N^{-1}(p)$. Given $Z$, the
default probability is $N\bigl((N^{-1}(p) - \sqrt\rho Z)/\sqrt{1-\rho}\bigr)$ and the
defaults are independent, so by the law of large numbers the defaulted fraction
equals it; the loss is $1-R$ times that.
\iqlookfor{conditional independence and the large-pool limit.}
\end{solution}

\begin{solution}{iq:m2:credit-indices-and-tranches:4}
Equity is long correlation (its expected loss falls as defaults cluster), senior
\omterm{def:m2:credit-indices-and-tranches:tranche}{tranches} short. The mezzanine's expected loss is not monotone in correlation, so one
price can fit two correlations, or none; base \omterm{def:m2:credit-indices-and-tranches:tranche}{tranches} all start at zero, and their
expected losses fall monotonically with correlation, so each price gives one
correlation.
\iqlookfor{direction of each \omterm{def:m2:credit-indices-and-tranches:tranche}{tranche}, and monotonicity.}
\end{solution}

\begin{solution}{iq:m2:credit-indices-and-tranches:5}
A portfolio meant as a hedge grew into a very large directional and relative-value book
without notional limits; the net notional tripled in a quarter; positions became a large
multiple of the market's daily volume, so they could not be exited without moving
prices; the marks were not independently controlled; and the losses were first
understated.
\iqlookfor{size against liquidity, limits, valuation control.}
\end{solution}

\begin{solution}{iq:m2:credit-indices-and-tranches:6}
Subscribe to single-name quotes and index quotes; keep each constituent's upfront and
annuity updated on each tick, with its weight in every series it belongs to; the
intrinsic upfront of a series is a weighted sum, updated incrementally when a name
moves; convert to a spread with a cached annuity and a Newton step; publish the skew
with the age and quality of the quotes behind it, and flag stale names.
\iqlookfor{incremental sums, many series per name, stale data.}
\end{solution}
